\section{Annexes}

\subsection{ACCOUSTIQUE}

\begin{table}[h!]
\centering
\begin{tabular}{|c|c|}
\hline
	$D_r_e_e_l$ \ cm &
	$\lambda$ \ cm  & \\
\hline
1,5 & 0,5 \\
2,0 & 0,9 \\
2,9 & 0,8 \\
3,7 & 1,8 \\
5,5 & 1,0 \\
6,5 & 0,9 \\
7,4 & 0,7 \\
8,1 & 0,9 \\
9,0 & N/A \\
\hline
\end{tabular}
\caption{Distances mesurées et nombre de longueurs d’onde correspondantes}
        \label{tab:lambda}
\end{table}

\subsection{ELECTRIQUE}
Entre 0 et 50 $\ohm$ :

\begin{table}[h!]
\centering
\begin{tabular}{|c|c|c|c|}
\hline
$R$ ($\Omega$) &
Amplitude Générateur (V) &
Amplitude Câble (V) &
Temps (ms) --> verifier \\
\hline
5   & 10.65 & 9.73 & 0.98 \\
10  & 9.65  & 7.48 & 1.79 \\
15  & 9.65  & 8.44 & 0.90 \\
20  & 8.84  & 6.43 & 1.83 \\
25  & 9.05  & 7.56 & 0.98 \\
30  & 8.19  & 5.95 & 1.42 \\
35  & 8.44  & 7.08 & 1.03 \\
40  & 8.04  & 5.39 & 1.59 \\
\hline
\end{tabular}
\caption{Mesures pour les résistances comprises entre 0 et 50 $\Omega$}
\end{table}

Pour R $>$ 50 $\ohm$ :

\begin{table}[h!]
\centering
\begin{tabular}{|c|c|c|c|}
\hline
$R$ ($\Omega$) &
Amplitude Générateur (V) &
Amplitude Câble (V) &
Temps (ns) \\
\hline
100  & 7.84 & 6.27 & 0.26 \\
500  & 8.24 & 8.68 & -1.87 \\
1500 & 7.64 & 9.33 & -2.51 \\
2000 & 8.04 & 9.49 & -2.59 \\
2500 & 7.84 & 9.43 & -2.79 \\
\hline
\end{tabular}
\caption{Mesures pour les résistances supérieures ou égales à 50 $\Omega$}
\end{table}

Synthese :

\begin{table}[h!]
\centering
\begin{tabular}{|c|c|c|c|}
\hline
$R$ ($\Omega$) &
Amplitude Générateur (V) &
Amplitude Câble (V) &
Temps (ns) \\
\hline
0     & N/A & N/A & N/A \\
50    & N/A & N/A & N/A \\
2500  & 7.84 & 9.43 & -2.79 \\
\hline
\end{tabular}
        \caption{Comparaison pour $R = 0\,\Omega$, $50\,\Omega$ et $R_{max}$ Ref \cite{christophe}}
\end{table}




\begin{thebibliography}{99}

\bibitem{rlc}
Cours electronique, \textit{Pont RLC},
\\
                \texttt{https://courselectronique.wordpress.com/etude-des-circuits/circuits-en-regime-transitoire/circuit-rlc-serie/}
\\
        Consulte le : 7 Janvier 2026

\end{thebibliography}

